\section*{Capítulo \ref{ch:b3:planetary-atmospheres} --- Atmósferas y habitabilidad planetaria}

\begin{solution}{exo:b3:planetary-atmospheres:1}
$H = RT/Mg$: Venus $8.314\times737/(0.04401\times8.87) = \qty{15.7}{km}$;
la Tierra $8.314\times288/(0.02897\times9.80) = \qty{8.4}{km}$; Marte
$8.314\times214/(0.04334\times3.71) = \qty{11.1}{km}$; Titán
$8.314\times94/(0.0280\times1.35) = \qty{20.7}{km}$. La atmósfera de Titán,
por fría que sea, es la más extensa: su gravedad es siete veces menor
que la de la Tierra.
\end{solution}

\begin{solution}{exo:b3:planetary-atmospheres:2}
$\lambda = GMm/kTr = \num{6.674e-11}\times\num{5.97e24}\times\num{1.674e-27}
/(\num{1.381e-23}\,T\times\num{6.878e6})$: 8.78 a \qty{800}{K}, 4.39 a
\qty{1600}{K}. Las fracciones de escape son $9.78\,e^{-8.78} =
\num{1.5e-3}$ y $5.39\,e^{-4.39} = \num{6.7e-2}$: el ciclo solar
modifica la fracción de escape del hidrógeno en un factor de unos 44.
\end{solution}

\begin{solution}{exo:b3:planetary-atmospheres:3}
(a)~$d = S_\mathrm{eff}^{-1/2}$: \omterm{def:b3:planetary-atmospheres:runaway}{invernadero húmedo} \qty{0.993}{au}, \omterm{def:b3:planetary-atmospheres:runaway}{límite del
invernadero desbocado} \qty{0.975}{au}, invernadero máximo \qty{1.705}{au}. (b)~Venus, en 1.91,
recibe más que el flujo \omterm{def:b3:planetary-atmospheres:runaway}{límite del invernadero desbocado} y aún más que el
límite empírico de Venus reciente (1.78): queda dentro del borde interior. Marte, en 0.43,
queda entre 1.014 y 0.3438, dentro de la zona habitable conservadora:
su sequedad procede de su pequeño tamaño y de la pérdida de su atmósfera, no
de su distancia.
\end{solution}

\begin{solution}{exo:b3:planetary-atmospheres:4}
(a)~$\log_{10}p = 3.9895 - 443.028/93.66 = -0.741$, $p =
\qty{0.18}{bar}$. (b)~$0.05\times1.6 = \qty{0.080}{bar}$, una
humedad relativa de $0.080/0.18 = 44\,\%$. (c)~El aire sobre los mares está lejos de
saturación: el metano se evapora desde ellos, como el agua de los
océanos de la Tierra, y el aire húmedo debe ascender y enfriarse antes de llover.
\end{solution}

\begin{solution}{exo:b3:planetary-atmospheres:5}
$\lambda_\mathrm D = 7.02\times2.014/1.008 = 14.0$. $u_\mathrm D/u_\mathrm H
= \sqrt{1.008/2.014} = 0.71$. Razón por átomo $= 0.71\times(15.0\,
e^{-14.0})/(8.02\,e^{-7.02}) = 0.71\times\num{1.69e-3} = \num{1.2e-3}$.
El \omterm{def:b3:planetary-atmospheres:jeans-escape}{escape de Jeans} deja salir casi nada de deuterio: es un filtro fuerte, con
$\alpha \approx 0.001$.
\end{solution}

\begin{solution}{exo:b3:planetary-atmospheres:6}
(a)~$0.016/\num{1.5575e-4} = 103$. (b)~Con $\alpha = 0$ el enriquecimiento es
$M_0/M$: $M_0 = 103\times\num{6e15} = \qty{6.2e17}{kg}$, o bien
$\num{6.2e17}/(1000\times\num{4.60e14}) = \qty{1.3}{m}$ de agua sobre
Venus. (c)~$M_0/M = 103^{1/0.2} = \num{1.2e10}$, un inventario pasado de
\qty{7e25}{kg}, doce veces la masa de la Tierra: imposible. El
enriquecimiento requiere un escape fuertemente fraccionador (de Jeans o no térmico),
activo al menos durante la pérdida de la última parte del agua, o un
estado estacionario con aporte de agua.
\end{solution}

\begin{solution}{exo:b3:planetary-atmospheres:7}
(a)~En cuatro mil millones de años ($\num{1.26e17}$~s) Marte pierde
$\num{3.8e43}$ átomos, $\num{1.9e43}$ moléculas de agua, $\qty{5.7e17}{kg}$;
sobre $4\pi(\num{3.3895e6})^2 = \qty{1.44e14}{m^2}$, lo que supone una capa
de \qty{3.9}{m} de espesor. (b)~Veinte veces más, \qty{78}{m}. (c)~Incluso la
tasa de alta oblicuidad alcanza solo el extremo inferior del intervalo de 100 a
\qty{1500}{m}; el resto pasó a la corteza y al hielo, o se perdió
más deprisa cuando el Sol era joven.
\end{solution}

\begin{solution}{exo:b3:planetary-atmospheres:8}
(a)~$m = pA/g = 636\times\num{1.44e14}/3.71 = \qty{2.5e16}{kg}$.
(b)~$0.1\times\num{1.26e17} = \qty{1.3e16}{kg}$, la mitad de la
atmósfera actual; toda ella en $\num{2.5e16}/0.1 = \qty{2.5e17}{s}$,
\num{7.8}~mil millones de años. (c)~El Sol joven era más activo y sus
tormentas eran más frecuentes, y la pérdida aumenta durante ellas; otros canales
(hidrógeno, pulverización catódica, escape fotoquímico) no se contabilizan.
\end{solution}

\begin{solution}{exo:b3:planetary-atmospheres:9}
(a)~$S(1 - 0.193)/4 = 291$ da $S = \qty{1442}{W.m^{-2}}$, 1.06 veces
la de la Tierra. (b)~Venus hoy absorbería $0.807\times2601/4 =
\qty{525}{W.m^{-2}}$; un Venus joven, $0.7\times2601 = \qty{1821}{W.m^{-2}}$
recibidos, \qty{367}{W.m^{-2}} absorbidos: ambos por encima del límite, así que un
Venus joven sin nubes habría entrado en \omterm{def:b3:planetary-atmospheres:runaway}{invernadero desbocado} desde el principio. (c)~Las nubes
elevan el albedo; los modelos tridimensionales con nubes han encontrado una vía
hacia un Venus joven habitable.
\end{solution}

\begin{solution}{exo:b3:planetary-atmospheres:10}
(a)~$H = kT/mg = \num{1.381e-23}\times1000/(16\times\num{1.661e-27}\times9.21)
= \qty{56}{km}$. (b)~$n_\mathrm{exo} = 1/(\sigma_c H) = 1/(\num{3e-19}\times
\num{5.64e4}) = \qty{5.9e13}{m^{-3}}$; $z = 200 + 56.4\ln(\num{1e16}/
\num{5.9e13}) = 200 + 56.4\times5.13 = \qty{490}{km}$. (c)~La gravedad, más débil
a mayor altura, hace que la altura de escala crezca, así que la densidad
disminuye más lentamente de lo supuesto: la \omterm{def:b3:planetary-atmospheres:exobase}{exobase} asciende.
\end{solution}

\begin{solution}{exo:b3:planetary-atmospheres:11}
(a)~$\Phi = \num{3.0e13}\times\num{3e-3} = \qty{9e10}{cm^{-2}.s^{-1}} =
\qty{9e14}{m^{-2}.s^{-1}}$; sobre $4\pi(\num{6.371e6})^2 =
\qty{5.10e14}{m^2}$, $\num{4.6e29}$ moléculas de $\ce{H2}$ por segundo.
(b)~Un océano contiene $\num{1.4e24}/18.015\times\num{6.022e23} =
\num{4.7e46}$ moléculas de agua, una de $\ce{H2}$ por cada una: dura
$\qty{1.0e17}{s} = \qty{3.2}{Gyr}$, comparable con la edad de
la Tierra, lo que explica que una razón de mezcla de \num{3e-3} marque el
límite del \omterm{def:b3:planetary-atmospheres:runaway}{invernadero húmedo}.
\end{solution}

\begin{solution}{exo:b3:planetary-atmospheres:12}
(a)~$p_1 = g/D\kappa = 9.81/(1.66\times0.005) = \qty{1180}{Pa}$.
(b)~$1/T = 1/273.16 - \ln(1182/611.657)\times461.5/\num{2.5e6}$, $T =
\qty{282.5}{K}$, y $\sigma T^4 = \qty{361}{W.m^{-2}}$. (c)~El límite real
es menor, así que el nivel radiante es más frío y está a menor presión: el
vapor de agua real es más opaco que este gas gris, y su $\kappa$ efectivo
es mayor.
\end{solution}

\begin{solution}{pb:b3:planetary-atmospheres:1}
\textbf{1.} $g = GM/R^2 = \num{6.674e-11}\times\num{5.2e24}/(\num{6.2e6})^2 =
\qty{9.03}{m.s^{-2}}$; $\bar\rho = M/\tfrac43\pi R^3 = \qty{5210}{kg.m^{-3}}$.

\textbf{2.} $v_\mathrm{esc} = \sqrt{2GM/r} = \sqrt{2\times\num{3.47e14}/\num{6.7e6}}
= \qty{10.18}{km.s^{-1}}$.

\textbf{3.} $u = \sqrt{2kT/m}$: hidrógeno \qty{4.81}{km.s^{-1}}, helio
\qty{2.41}{km.s^{-1}}, oxígeno \qty{1.21}{km.s^{-1}}.

\textbf{4.} $\lambda = (v_\mathrm{esc}/u)^2$: 4.49, 17.8, 71.2.

\textbf{5.} Solo el hidrógeno: su fracción de escape es 0.06, la del helio
\num{3.5e-7}, la del oxígeno \num{9e-30}. Ninguno está cerca del régimen hidrodinámico,
$\lambda \approx 1.5$ a 3.

\textbf{6.} $\lambda = 4.49\times1400/3000 = 2.09$: dentro de la transición;
el hidrógeno habría escapado hidrodinámicamente, quizá arrastrando
gases más pesados.

\textbf{7.} $(1 + 4.49)\,e^{-4.49} = 0.0618$.

\textbf{8.} $n u/2\sqrt\pi = \num{2e11}\times4806/3.545 = \qty{2.71e14}{m^{-2}.s^{-1}}$;
$\Phi_\mathrm J = \num{2.71e14}\times0.0618 = \qty{1.68e13}{m^{-2}.s^{-1}}$.

\textbf{9.} $4\pi(\num{6.7e6})^2 = \qty{5.64e14}{m^2}$; $\num{9.5e27}$ átomos por
segundo.

\textbf{10.} $\num{1.4e24}/18.015 = \qty{7.8e22}{mol}$, $\num{4.68e46}$ moléculas,
$\num{9.36e46}$ átomos de hidrógeno.

\textbf{11.} $\num{9.36e46}/\num{9.5e27} = \qty{9.9e18}{s}$, unos 310 mil millones de
años.

\textbf{12.} $\lambda_\mathrm D = 4.49\times2.014/1.008 = 8.96$; fracción de escape
$9.96\,e^{-8.96} = \num{1.28e-3}$; eficiencia relativa $0.707\times
\num{1.28e-3}/0.0618 = 0.015$.

\textbf{13.} $H_a = kT/m_a g = \num{1.381e-23}\times200/(44.01\times\num{1.661e-27}
\times9.03) = \qty{4.18}{km}$.

\textbf{14.} $b = \qty{1.2e21}{m^{-1}.s^{-1}}$, así que
\[
\Phi_\mathrm{DL} = \frac{b f(m_a - m_i)}{m_a H_a}
= \frac{\num{1.2e21}\times\num{1e-3}\times(42.0/44.0)}{4185}
= \qty{2.74e14}{m^{-2}.s^{-1}}.
\]

\textbf{15.} $\num{3.0e13}\times\num{1e-3}\times9.03/9.81 =
\qty{2.76e10}{cm^{-2}.s^{-1}} = \qty{2.76e14}{m^{-2}.s^{-1}}$: lo mismo
dentro del 1\,\%.

\textbf{16.} $4\pi(\num{6.2e6})^2 = \qty{4.83e14}{m^2}$; $\num{1.32e29}$ moléculas
de $\ce{H2}$, $\num{2.64e29}$ átomos por segundo.

\textbf{17.} El límite de difusión aporta 28 veces lo que la \omterm{def:b3:planetary-atmospheres:exobase}{exobase} deja escapar con
su densidad actual. El hidrógeno se acumula en la termosfera hasta que el
flujo de Jeans iguale al aporte: a largo plazo el escape está
limitado por difusión, como en la Tierra.

\textbf{18.} $\num{2e11}\times28 = \qty{5.6e12}{m^{-3}}$.

\textbf{19.} $\num{4.68e46}/\num{1.32e29} = \qty{3.54e17}{s}$, 11.2 mil millones de años.

\textbf{20.} $S = 0.30/0.60^2 = 0.833$.

\textbf{21.} $d = \sqrt{L/S_\mathrm{eff}}$: borde interior $\sqrt{0.30/0.939} =
\qty{0.565}{au}$, borde exterior $\sqrt{0.30/0.290} = \qty{1.02}{au}$. Kepos,
a \qty{0.60}{au}, está dentro, cerca del borde interior.

\textbf{22.} $0.939/0.833 = 1.13$: un aumento de la insolación del 13\,\%.

\textbf{23.} $0.833\times1361\times0.7/4 = \qty{198}{W.m^{-2}}$, $T_\mathrm{eq} =
(198/\sigma)^{1/4} = \qty{243}{K}$; muy por debajo del
límite de \qty{291}{W.m^{-2}} (calculado para una estrella similar al Sol, pero el margen
es amplio).

\textbf{24.} Tres veces más rápido: 3.7 mil millones de años.

\textbf{25.} \textbf{11~mil~millones~de~años}.
\end{solution}
